% Restored and adapted detailed content from the previous introduction, background, and related-work draft.

% Source: paper/chapters/01-introduction.tex

Music composition is an inherently structural discipline.
Classical sonata forms, jazz chord progressions,
and electronic dance music loops all rely on the hierarchical organization of motifs, phrases, and sections.
Despite this underlying mathematical and logical foundation, modern digital music production relies
predominantly on linear, visual interfaces.
While these tools excel at micro-editing, they often lack the abstractions necessary to manipulate music as structured,
algorithmic data.

\section{Motivation}\label{detail-sec:motivation}

In a standard Digital Audio Workstation (DAW), the primary method of input is the ``Piano Roll''---a grid
where notes are
drawn manually across a timeline.
This approach is highly intuitive but fundamentally destructive from a structural perspective.
If a composer wishes to transpose a repeating 16-bar melody, change its rhythm, or conditionally alter its ending on the
final repetition, they must manually select and modify hundreds of individual note events.
The logical intent (e.g., ``repeat this phrase four times but play the snare on the off-beat during the last
repetition'') is lost, replaced by a static collection of data points.

To address this, programmers often turn to Algorithmic Composition.
However, existing domain-specific languages for music (such as Sonic Pi or Tidal) are largely tailored toward
\textit{live coding}.
They prioritize real-time performance, utilizing threaded execution models and continuous audio synthesis.
While incredibly powerful for improvisation, these environments introduce complexities related to thread
synchronization, execution latency, and audio engine configuration, which can be overwhelming for composers whose
primary goal is offline, deterministic arrangement.

Furthermore, general-purpose programming languages equipped with MIDI libraries (e.g., Python with \texttt{mido}) offer
offline generation but lack domain-specific abstractions.
The composer is burdened with calculating explicit delta-times and manually constructing low-level binary messages,
shifting the focus away from musical creativity and toward byte-level engineering.

\section{Problem Statement}\label{detail-sec:problem-statement}

There exists a gap in the current ecosystem of music programming tools: a lack of a domain-specific language optimized
for \textit{offline, structural composition} that produces standard, interoperable artifacts.
Specifically, the challenge lies in creating a system that allows composers to define musical ideas as reusable,
parameterized algorithms while keeping the resulting temporal structure deterministic, free of real-time
execution jitter,
and readable by many standard production tools.

\section{Proposed Solution}\label{detail-sec:proposed-solution}

This thesis proposes the design and implementation of a novel, compiled Domain-Specific Language
(domain-specific language) for musical
composition.
The system aims to bridge the gap between algorithmic logic and the rigid binary specifications of the MIDI protocol.

The proposed solution is built upon the following core principles:

\noindent \textbf{Music as a Recursive Hierarchy.} The language abstracts musical motifs into functions (patterns) that
can be parameterized, reused, and nested, reflecting the true hierarchical nature of musical form.

\noindent \textbf{C-Style Imperative Flow.} By utilizing familiar control flow structures (\texttt{for}, \texttt{loop},
\texttt{if}), Motivo allows composers to procedurally generate variations and repetitions without learning esoteric
functional paradigms.

\noindent \textbf{Zero-Runtime Compilation.} The compiler operates strictly offline.
It performs semantic validation and then uses a stateful lowering engine to completely unroll the composition's
logic into a flat timeline.
This keeps temporal resolution deterministic within the compiler pipeline without relying on an interpreter
during playback.

\noindent \textbf{MIDI as a Semantic Target.} The compiler serializes the final timeline into a Standard MIDI File
(Format 1).
This supports practical interoperability, allowing the generated compositions to be imported into many DAWs
and MIDI-aware tools for
professional orchestration and mixing.

% Source: paper/chapters/02-background/music_theory.tex

\section{Music Theory Primer}\label{detail-sec:background-music-theory}

To understand the abstractions provided by a musical domain-specific language, it is necessary to establish a
foundational understanding of western music theory.
Music is organized both vertically (pitch) and horizontally (time).

\subsection{Pitch, Notes, and Octaves}\label{detail-subsec:pitch-notes-and-octaves}

The fundamental unit of melody and harmony is the \textit{note}, which represents a sound of a specific frequency.
In western music, the octave is divided into twelve semitones.
The seven natural notes are designated by the letters A through G\@.
\textit{Accidentals} modify these natural notes: a sharp (\#) raises the pitch by one semitone, while a flat (b) lowers
it by one semitone.

An \textit{octave} is the interval between one musical pitch and another with double its frequency.
To uniquely identify a note across the entire audible spectrum, it is denoted by its pitch class and its octave number
(e.g., \texttt{C4}, which is often referred to as ``Middle C'').
When two or more distinct notes are sounded simultaneously, they form a \textit{chord}, which provides the harmonic
foundation of a composition.

\subsection{Time, Rhythm, and Tempo}\label{detail-subsec:time-rhythm-and-tempo}

While pitch defines the frequency of a sound, rhythm defines its placement and duration in time.
The speed at which a piece of music is played is governed by its \textit{tempo}, typically measured in Beats Per Minute
(BPM).
A tempo of 120 BPM signifies that there are 120 quarter-note beats in one minute, or one beat every 0.5 seconds.

Music is structurally divided into \textit{measures} (or bars) by a \textit{time signature}.
A time signature, such as 4/4, consists of two numbers: the numerator indicates the number of beats per measure (four),
and the denominator indicates the note value that constitutes one beat (the quarter note).
A \textit{rest} is a period of silence that has a specific duration, functioning as a ``silent note'' that preserves the
rhythmic alignment of the timeline.

% Source: paper/chapters/02-background/midi_protocol.tex

\section{The MIDI Protocol}\label{detail-sec:background-midi}

The Musical Instrument Digital Interface (MIDI) is an industry-standard technical protocol developed in 1983 to enable
electronic instruments, computers, and other related devices to communicate and synchronize with each other.
Crucially, MIDI does not transmit digital audio (sound waves); instead, it transmits \textit{event messages} that act as
instructions for how music should be produced by a receiving synthesizer.

\subsection{Event-Based Representation}\label{detail-subsec:event-based-representation}

A MIDI stream is composed of sequential commands.
The most fundamental commands are:

\noindent \textbf{Note-On.} Instructs the synthesizer to begin sounding a pitch.
It carries two data bytes: the \textit{key number} (0--127, where 60 is typically Middle C) and the
\textit{velocity} (0--127, representing how hard the key was struck).

\noindent \textbf{Note-Off.} Instructs the synthesizer to stop sounding the pitch.

\noindent \textbf{Program Change.} Changes the active instrument (or ``patch'') on a specific channel (e.g., switching
from a piano to a violin).

Because MIDI transmits events rather than waveforms, a piece of music represented in MIDI is highly editable.
A composer can retroactively change the pitch, duration, or instrument of any note without degrading the audio quality,
making MIDI the \textit{de facto} standard for digital composition.

\subsection{Channels and General MIDI}\label{detail-subsec:channels-and-general-midi}

To support polyphony across different instruments, MIDI data is transmitted across 16 independent channels.
A synthesizer can be configured to play a bass patch on Channel 1 and a piano patch on Channel 2 simultaneously.
The \textit{General MIDI} (GM) specification standardizes the assignment of instruments to program numbers (e.g.,
Program 1 is Acoustic Grand Piano) and strictly reserves \textbf{Channel 10} for unpitched percussion, where individual
key numbers trigger different drum samples (e.g., Key 36 is the Bass Drum).

\subsection{Delta-Times and Standard MIDI Files (SMF)}\label{detail-subsec:delta-times-and-standard-midi-files-(smf)}

When saving MIDI data to a file (a \texttt{.mid} file), the protocol utilizes \textit{delta-times}.
Instead of recording the absolute time at which an event occurs, a Standard MIDI File records the number of ``ticks'' to
wait after the \textit{previous} event before executing the current one.
Ticks are subdivisions of a quarter note, defined by the file's \textit{Pulses Per Quarter-note} (PPQ) resolution (e.g.,
480 PPQ).
This relative timing system is highly efficient but requires careful calculation when translating from absolute musical
beats into a serialized file format.

% Source: paper/chapters/02-background/compiler_basics.tex

\section{Compiler Construction Basics}\label{detail-sec:background-compilers}

A compiler is a specialized software system that translates source code written in a high-level programming language
into a lower-level representation, such as machine code, bytecode, or in the context of this thesis, a binary MIDI file.
The compilation process is traditionally divided into several distinct phases to manage complexity and ensure
correctness.

\subsection{Lexical and Syntactic Analysis}\label{detail-subsec:lexical-and-syntactic-analysis}

The first phase, \textit{Lexical Analysis} (or scanning), reads the raw character stream of the source file and groups
characters into meaningful sequences called \textit{tokens}.
For example, the string \texttt{tempo} might be converted into a \texttt{KEYWORD\_TEMPO} token, while \texttt{120}
becomes an \texttt{INT\_LITERAL} token.
This is often implemented using tools like Flex, which rely on regular expressions.

The second phase, \textit{Syntactic Analysis} (or parsing), takes the stream of tokens and verifies that they form a
valid sequence according to the formal grammar of the language.
This phase organizes the flat token stream into a hierarchical tree structure known as an \textit{Abstract Syntax Tree}
(AST).
The AST captures the logical operations of the code (e.g., nesting a \texttt{play} statement inside a \texttt{loop}
block) while discarding irrelevant textual details like whitespace and semicolons.

\subsection{Semantic Analysis}\label{detail-subsec:semantic-analysis}

While the parser ensures that a sentence is grammatically correct, the \textit{Semantic Analyzer} ensures that it makes
logical sense.
This phase traverses the AST to perform checks that cannot be captured by a context-free grammar.
Key tasks include:

\noindent \textbf{Type Checking.} Ensuring that operations are applied to compatible data types (e.g., preventing the
addition of a boolean to a note).

\noindent \textbf{Symbol Resolution.} Maintaining a \textit{Symbol Table} to ensure that variables and functions are
declared before they are used, and that they are accessed within their valid scope.

\subsection{Intermediate Representation and Code
Generation}\label{detail-subsec:intermediate-representation-and-code-generation}

After validation, the compiler often translates the AST into an \textit{Intermediate Representation} (IR).
An IR is a lower-level, often flattened representation of the program that is easier to optimize and translate than the
hierarchical AST\@.
In standard compilers, this might involve generating three-address code.
In the context of music generation, lowering to an IR involves unrolling loops and resolving relative timings into
absolute timestamps.

Finally, the \textit{Code Generation} phase translates the IR into the target output format.
This phase handles the intricate binary serialization details, such as little/big-endian encoding and specific protocol
constraints, yielding the final executable artifact.

% Source: paper/chapters/02-background/algorithmic_composition.tex

\section{Algorithmic Composition}\label{detail-sec:background-algorithmic-composition}

Algorithmic composition is the technique of using algorithms---formal sets of rules or instructions---to create music.
While this concept has roots in classical music (such as Mozart's musical dice games), the advent of computers has
exponentially expanded its scope and capabilities.

\subsection{Procedural Generation vs. Direct Input}\label{detail-subsec:procedural-generation-vs.-direct-input}

Traditional digital music production relies on Direct Input.
A composer uses a MIDI keyboard or a graphical Piano Roll in a Digital Audio Workstation (DAW) to manually input every
note, chord, and rest.
This process grants total control over the micro-details of a performance but can be tedious when constructing complex,
repetitive, or mathematically structured sections.

Algorithmic composition shifts the paradigm to Procedural Generation.
Instead of writing the notes, the composer writes the \textit{rules} that generate the notes.
For example, rather than manually copying and pasting a bassline across 32 bars while applying slight variations to the
velocity, a composer might write a simple loop that algorithmically modifies the velocity parameter on each iteration.

\subsection{Structured Abstraction}\label{detail-subsec:structured-abstraction}

The true power of algorithmic composition lies in its ability to abstract musical ideas.
A chord progression can be represented as an array of data points, and a rhythm can be defined as a mathematical
function over time.
This approach allows composers to experiment rapidly: changing a foundational variable (such as a root note or a time
signature) can instantly cascade throughout the entire algorithmic structure, generating a completely new variation of
the piece without the need for manual editing.
Motivo proposed in this thesis aims to provide a robust, type-safe environment specifically optimized for this
methodology of structured, algorithmic creation.

% Source: paper/chapters/03-related-work/introduction.tex

\section{Introduction}\label{detail-sec:introduction}

This chapter surveys existing approaches to music representation and programming, focusing on systems that enable the
creation, manipulation, and execution of musical structures through textual or programmatic means.

The reviewed works span multiple categories, including high-level composition systems, low-level audio programming
languages, and music representation formats.
These systems differ significantly in their design goals, levels of abstraction, and approaches to modeling musical
time and structure.

High-level composition systems, such as \textbf{Sonic Pi}, \textbf{Tidal}, and \textbf{HMusic}, emphasize expressive
musical abstractions, often prioritizing accessibility, structured composition, and interactive use.

In contrast, low-level audio programming languages like \textbf{ChucK} and \textbf{SuperCollider} provide fine-grained
control over sound synthesis and timing, but operate at a lower level of abstraction.

Additionally, music representation formats such as \textbf{MIDI} and \textbf{ABC notation} define how musical data can
be encoded, stored, and transmitted, often prioritizing interoperability and readability over compositional abstraction.

By analyzing these systems, this chapter identifies key design trade-offs and limitations in existing approaches,
particularly with respect to temporal modeling, abstraction level, and compositional structure.
These observations serve as the foundation for motivating the design of the domain-specific language proposed in
this thesis.

% Source: paper/chapters/03-related-work/high-level-musical-composition/sonic_pi.tex

\subsection{Sonic Pi}\label{detail-subsec:sonic-pi}

\textbf{Sonic Pi} began as a 3-month pilot project in \textit{November 2013}, funded by the \textit{Broadcom
Foundation}.
It aimed to teach school students with no prior programming experience core computational concepts through a live coding
environment, as well as enable music composers to deliver live coded performances~\cite{Aaron_Blackwell_2013,
Aaron_2016}.

This language was developed to run on the \textit{Raspberry Pi}, having to work on tight hardware constraints, where
time and memory overhead needed to be managed carefully.
It's core architecture was designed around an earlier project called \textit{Overtone}, but which had too steep of a
learning curve for music performers and had massive performance issues on the \textit{Raspberry Pi}~\cite{Aaron_2016}.

%% ---------------------------------------------------------------------------------------------------------------------

\subsubsection{Key Features}

To address the tight constraints of having to be optimized to run on low end hardware such as the \textit{Raspberry Pi},
the language was designed and embedded in \textit{Ruby}.

The system is designed in terms of the \textit{SuperCollider} sound synthesis engine, allowing the control
and manipulation of synthesizers in real time.
\textit{SuperCollider} acts as the core music synthesis backend for \textbf{Sonic Pi}, as it has no synthesis
capabilities of its own~\cite{Aaron_Blackwell_2013}.

\textbf{Sonic Pi} follows a \textit{sequential execution model}, where musical events are generated in the order in
which instructions are executed.
We can see this in action in the following code snippets explaining its core methods.

The simplest way to play a sound is through the \texttt{play} method, which takes a \textit{MIDI} note as its argument.
Notes can be separated and spaced out between each other using the \texttt{sleep} method, which takes a floating point
number as its argument, representing the time in seconds to wait until playing the next sound.
Notes can be combined into patterns by using simple lists, and they can be played sequentially using the
\texttt{play\_pattern} method~\cite{Aaron_Blackwell_2013}.

Example snippet:
\begin{lstlisting}[label={detail-lst:sonic-pi-1}]
  play 60;
  sleep 0.5;
  play_pattern [50, 55, 60];
\end{lstlisting}

\textbf{Sonic Pi} also introduces simple control flow structures such as \textit{loops} and
\textit{conditional statements}:
\begin{lstlisting}[label={detail-lst:loops-and-conditional-statements}]
  10.times do
    play 65;
    sleep 0.25;
    ...
  end

  if num < 0.5
    play 60
  else
    play 70
  end
\end{lstlisting}

One more powerful feature is the ability to stack options or filters on a certain node, by specifying the name of the
option and a value, enabling complex sound modulation and more granularity and control over notes.
Examples of these options include \texttt{attack}, \texttt{release}, \texttt{sustain} and many more.

This simplistic model enables an easier grasp over musical compositional concepts for users with non-technical
backgrounds, and is the main reason why it had such a positive impact in teaching sessions in middle schools in the
UK~\cite{Aaron_2016, Aaron_Blackwell_2013}.

%% ---------------------------------------------------------------------------------------------------------------------

\subsubsection{Limitations}

\textbf{Sonic Pi} functions more as an \textit{OSC} client that communicates to a subset of the \textit{SuperCollider}
API\@.
This subset, although small, is just enough for triggering sounds, controlling synthesizers and clearing data between
runs.
This, however, leaves \textbf{Sonic Pi} incapable of synthesizing sound on its own, and thus relies on an external
synthesis engine~\cite{Aaron_Blackwell_2013}.

Because it is primarily an educational tool with a sequential execution model, Sonic Pi offers limited compositional
scalability and structural abstraction for the offline structural goals of this thesis.
Complex layering capabilities are sacrificed for better ease of use and mental models, leaving the language weaker than
\textit{Overtone}, on which it was based.

%% ---------------------------------------------------------------------------------------------------------------------

\subsubsection{Relevance to Motivo}

The syntax of this language is both clean and concise, thus making it easier to map abstract concepts from musical
theory into textual form.
It creates the building blocks necessary to construct more complex structures, providing a good basis for Motivo's key
syntactical model.
However, the sequential execution model can make it more difficult to explicitly represent complex temporal
relationships between independent musical structures, which motivates the timeline-based approach proposed in this
thesis.

% Source: paper/chapters/03-related-work/high-level-musical-composition/tidal.tex

\subsection{Tidal}\label{detail-subsec:tidal}

\textbf{Tidal} is a domain specific programming language embedded in \textit{Haskell} for encoding musical patterns,
meant to be used during live performances~\cite{McLean_Wiggins_2010}.
It adopts a pattern-based model of music representation, where musical structures are defined as transformations of
time-varying patterns rather than explicit event placement.
It is meant to address the gap in pattern languages, as most languages focused on encoding and decoding musical patterns
at that time were focused more on \textit{musical analysis} rather than \textit{musical
composition}~\cite{McLean_Wiggins_2010}.

%% ---------------------------------------------------------------------------------------------------------------------

\subsubsection{Key Features}

Since it is embedded in the programming language \textit{Haskell}, it inherits most of its syntax and paradigms.
It uses Haskell's strict type system, which in term leads to stricter constraints such as static typing and pure
function usage.
However, it makes use of these constraints to its advantage by leveraging \textit{Haskell's} concepts of
\textit{functors} and \textit{monads}~\cite{McLean_Wiggins_2010}.

\paragraph{Pattern Composition}

Patterns can be composed of sub-patterns; this relation can go on to an arbitrary depth.
This allows for creating complex structures from simple building blocks, creating an organized, \textit{hierarchical}
structure.
Patterns can be concatenated, merged, or even combined by multiplying their numerical patterns, offering great control
and extensibility through pair-wise operations~\cite{McLean_Wiggins_2010}.

\paragraph{Representation and Access}

Some languages represent patterns by using lazily loaded lists, where values are calculated only when they are needed,
which allows for very long or even infinite patterns to be computed and represented.
\textit{Haskell} itself uses lazy lists by default, but \textbf{Tidal} decides to not follow the same idea.

Lazy lists don't work great for live performances since some elements are calculated based on previous items in that
list.
If you modify the definition of a lazy list, in order to continue from a set point, you'd have to recompute
the entire list up to that point, which is computationally expensive, especially in a live coding environment.
This is why \textbf{Tidal} allows \textit{random access} by use of \textit{functions} mapping time values to
events~\cite{McLean_Wiggins_2010}.

\paragraph{Timing}

Rhythmic structures progress linearly, thus it is not that complex to encode these progressions and allow for periodic
or infinite patterns, concepts that \textbf{Tidal} allows.
Timing can be either \textit{discrete} or \textit{continuous}, and musical tradition combines these concepts by often
using discrete notations (like the Western musical staff), paired with playing phrases subtly over continuous time.
\textbf{Tidal} reflects this by defining patterns as \textit{events over discrete time steps}, with the possibility of
defining some floating point onsets as well~\cite{McLean_Wiggins_2010}.

%% ---------------------------------------------------------------------------------------------------------------------

\subsubsection{Limitations}

That said, \textbf{Tidal} also comes with some limitations, such as its reliance on \textit{Haskell's} paradigms and
structure, sacrificing syntax for a higher learning curve, but with the benefit of powerful and efficient pattern
construction.

It also lacks compiling capabilities, as it is meant for outputting events and not actually synthesize sound on its own.
Thus, it relies on \textbf{OSC} (\textit{Open Sound Control}) to stream musical events to an external synthesizer.
This is of course due to its main target being live performances, but this in term makes music composition
for bedroom-producers much harder, as they would have to rely on external hardware~\cite{McLean_Wiggins_2010}.

%% ---------------------------------------------------------------------------------------------------------------------

\subsubsection{Relevance to Motivo}

\textbf{Tidal} describes how to encode one of the most fundamental constructs in music theory and musical composition.
It shows how we can make pattern declarations more consistent and concise through definitions, rather than hand-placed
notes.

These limitations, however, highlight a different design focus compared to Motivo proposed in this thesis,
which emphasizes deterministic composition and direct compilation to output formats.

% Source: paper/chapters/03-related-work/high-level-musical-composition/hmusic.tex

\subsection{HMusic}\label{detail-subsec:hmusic}

\textbf{HMusic} is a domain-specific language for music programming and live coding, embedded in the Haskell functional
programming language, and designed around the abstractions of patterns and tracks~\cite{Du_Bois_Ribeiro_2019}.

The language aims to bridge the gap between traditional music production tools, such as sequencers and drum machines,
and formal programming languages by providing a representation of music that resembles grid-based composition interfaces
while retaining the expressive power of functional programming~\cite{Du_Bois_Ribeiro_2019}.
Programs written in \textbf{HMusic} resemble the visual structure of multi-track sequencers, where patterns define
temporal structures and tracks associate these patterns with instruments, enabling both intuitive composition and
programmatic manipulation of musical structures~\cite{Du_Bois_Ribeiro_2019}.

%% ---------------------------------------------------------------------------------------------------------------------

\subsubsection{Key Features}

\textbf{HMusic} is centered around two main abstractions: \textit{patterns} and \textit{tracks}.
Patterns represent sequences of musical events, typically modeled as hits and rests, and are defined using
inductive data types that allow recursive composition.
Tracks associate patterns with instruments, enabling the separation of musical structure from sound generation.
Multiple tracks can be combined in parallel to form multi-track compositions, while sequential composition allows
the construction of longer musical structures from smaller components~\cite{Du_Bois_Ribeiro_2019}.

Example snippet:

\begin{lstlisting}[label={detail-lst:hmusic-1}]
  kick = X :| O :| O :| O
  snare = O :| O :| X :| O
  track = MakeTrack "BassDrum" kick
    :|| MakeTrack "Snare" snare
\end{lstlisting}

As the language is embedded in Haskell, it benefits from functional programming features such as higher-order functions
and recursion, allowing users to define reusable transformations over patterns and tracks.
\textbf{HMusic} also supports live coding through primitives for playing, looping, and modifying tracks in real time,
enabling interactive and dynamic composition workflows.

Additionally, the language is compiled into \textit{Sonic Pi}, leveraging its synthesis capabilities while focusing
on higher-level musical abstractions~\cite{Du_Bois_Ribeiro_2019}.

%% ---------------------------------------------------------------------------------------------------------------------

\subsubsection{Limitations}

Despite its expressive abstractions, \textbf{HMusic} presents several limitations.

First, the language depends on \textit{Sonic Pi} as a backend for sound synthesis, limiting its independence and
requiring integration with external tools for execution.

Second, as an embedded domain-specific language in Haskell, it inherits the complexity of functional
programming, which may present a barrier
for users without prior programming experience.

Furthermore, the pattern-based model primarily focuses on rhythmic and grid-like structures, which may limit its ability
to represent more complex or hierarchical musical compositions.

Finally, while patterns and tracks provide compositional flexibility, the language does not explicitly model
higher-level temporal structures such as timelines or sections, which can make large-scale composition more difficult to
manage.

%% ---------------------------------------------------------------------------------------------------------------------

\subsubsection{Relevance to Motivo}

\textbf{HMusic} demonstrates an effective approach to combining intuitive musical representations with formal
programming abstractions, particularly through its use of patterns and tracks as core compositional units.
Its design highlights the benefits of separating musical structure from sound generation and enabling composition
through reusable abstractions.

However, its reliance on pattern-based composition and external compilation to Sonic Pi contrasts with the goals of
this thesis, which aim to provide a more explicit and flexible temporal model.
Motivo builds upon these ideas by introducing a timeline-based approach, improved modularity, and direct
compilation to formats such as MIDI or audio, enabling more structured and deterministic composition workflows.

% Source: paper/chapters/03-related-work/real-time-audio-control-and-synthesis/chuck.tex

\subsection{ChucK}\label{detail-subsec:chuck}

\textbf{ChucK} is a strongly-timed, concurrent programming language for real-time audio synthesis, composition, and
performance, developed by Ge Wang and Perry R. Cook~\cite{Wang_Cook_2003}.
It was designed to provide precise control over time, enabling composers to express complex audio processes with
deterministic timing and concurrency.

The language introduces a model where time is treated as a first-class concept, allowing programmers to explicitly
control the progression of time and the execution of musical events.
This enables accurate synchronization of concurrent processes and fine-grained control over temporal relationships in
music~\cite{Wang_Cook_2003}.

%% ---------------------------------------------------------------------------------------------------------------------

\subsubsection{Key Features}

\textbf{ChucK} follows a \textit{strongly-timed execution model}, where time advances only when explicitly specified by
the programmer.
The special keyword \texttt{now} represents the current time, and execution remains suspended until time is advanced,
ensuring deterministic behavior across runs~\cite{Wang_Cook_2003}.

A central feature of the language is its support for \textit{concurrency} through lightweight processes called
\textit{shreds}.
These shreds execute in parallel and are synchronized through the same timing mechanism, allowing multiple musical
processes to run concurrently with precise temporal alignment.

Example snippet:

\begin{lstlisting}[label={detail-lst:chuck-1}]
  SinOsc osc => dac;
  440 => osc.freq;
  while (true) {
    1::second => now;
  }
\end{lstlisting}

Another defining characteristic is the \textit{ChucK operator} (\texttt{=>}), which unifies multiple language constructs
such as assignment, data flow, and signal routing.
This operator enforces a left-to-right execution model and allows chaining of operations, making the flow of audio and
data explicit~\cite{Wang_Cook_2003}.

The language also supports \textit{on-the-fly programming}, enabling code to be added, modified, or removed during
execution without interrupting the running system.
This feature makes \textbf{ChucK} particularly suitable for live coding and interactive performance
contexts~\cite{Wang_Cook_2003}.

Additionally, \textbf{ChucK} supports multiple simultaneous control rates, allowing different concurrent processes to
operate at varying temporal resolutions.

%% ---------------------------------------------------------------------------------------------------------------------

\subsubsection{Limitations}

Despite its precision and flexibility, \textbf{ChucK} presents several challenges.

First, the language operates at a relatively low level of abstraction, requiring programmers to explicitly manage time
and concurrency.
This can make it difficult to represent higher-level musical structures such as patterns, timelines, or hierarchical
compositions.

Second, the strongly-timed model places the responsibility of advancing time on the programmer, which can increase
complexity and reduce readability in larger programs.

Furthermore, while \textbf{ChucK} excels in real-time synthesis and performance, it is less focused on structured,
deterministic composition workflows, particularly those that require clear separation between composition and execution.

%% ---------------------------------------------------------------------------------------------------------------------

\subsubsection{Relevance to Motivo}

\textbf{ChucK} highlights the importance of explicit time control and deterministic concurrency in music programming,
demonstrating how precise temporal semantics can be integrated directly into a programming language.

However, its low-level and time-driven approach contrasts with the goals of this thesis, which aim to provide a
higher-level, structured representation of musical concepts.
Motivo seeks to retain deterministic behavior while offering concise abstractions for composition, such
as timeline-based structures and modular musical components.

% Source: paper/chapters/03-related-work/real-time-audio-control-and-synthesis/super_collider.tex

\subsection{SuperCollider}\label{detail-subsec:supercollider}

\textbf{SuperCollider} is a programming language and environment for real-time audio synthesis and algorithmic
composition, originally developed by James McCartney.
It was designed to provide a highly expressive and flexible platform for both sound synthesis and musical structure
definition, combining a high-level language with a powerful synthesis engine~\cite{McCartney_2002}.

Unlike traditional music programming environments, \textbf{SuperCollider} was built with the goal of supporting dynamic
and generative compositions, where musical processes can evolve over time rather than representing a fixed score.
It enables composers to describe not only concrete musical events, but also systems that generate music
algorithmically~\cite{McCartney_2002}.

%% ---------------------------------------------------------------------------------------------------------------------

\subsubsection{Key Features}

\textbf{SuperCollider} follows an \textit{object-oriented and functional programming model}, where all
elements are treated as objects and can be composed using high-level abstractions.
The language provides a rich set of programming constructs, including functions, control structures, and data types,
allowing for complex algorithmic manipulation of sound~\cite{McCartney_2002}.

A central abstraction in \textbf{SuperCollider} is the concept of a \textit{unit generator} (\textit{UGen}),
which represents a basic signal processing component.
These UGens can be combined to form synthesis graphs, defining how audio signals are generated and processed.
Instruments are constructed functionally, meaning that a program describes how to build a network of UGens rather than
specifying a fixed structure~\cite{McCartney_2002}.

Example snippet:

\begin{lstlisting}[label={detail-lst:supercollider-1}]
  {
    SinOsc.ar(440, 0, 0.2)
  }.play;
\end{lstlisting}

\textbf{SuperCollider} also introduces higher-level abstractions such as \textit{patterns} and \textit{streams}, which
allow composers to define sequences of musical events.
A stream represents a sequence of values generated over time, while patterns define reusable structures that can produce
multiple streams~\cite{McCartney_2002}.

Another important architectural feature is the separation between the \textit{language} and the
\textit{synthesis server}.
The language is responsible for generating musical events, while the synthesis engine executes them.
Communication between the two is handled through the \textit{Open Sound Control (OSC)} protocol, enabling distributed
and real-time control over audio processes~\cite{McCartney_2002}.

This architecture allows \textbf{SuperCollider} to support:

\noindent real-time audio synthesis

\noindent dynamic creation and modification of sound processes

\noindent distributed and network-based performance setups

%% ---------------------------------------------------------------------------------------------------------------------

\subsubsection{Limitations}

Despite its expressive power, \textbf{SuperCollider} presents several challenges.

First, the language has a \textit{steep learning curve}, largely due to its combination of object-oriented, functional,
and domain-specific paradigms.
Understanding concepts such as UGens, synthesis graphs, and event streams requires both programming knowledge and
familiarity with digital signal processing.

Second, the separation between the language and the synthesis server introduces additional complexity.
Communication via \textit{OSC} can lead to latency and requires users to reason about asynchronous execution
and message passing~\cite{McCartney_2002}.

Furthermore, while \textbf{SuperCollider} excels at sound synthesis and real-time performance, it is less focused on
structured, high-level composition.
Musical ideas are often expressed through low-level signal processing constructs, which can make it harder to represent
abstract musical structures such as timelines or hierarchical compositions.

%% ---------------------------------------------------------------------------------------------------------------------

\subsubsection{Relevance to Motivo}

\textbf{SuperCollider} demonstrates a powerful approach to integrating sound synthesis with a high-level
programming language, offering extensive flexibility in defining and manipulating audio processes.
Its use of composable abstractions, such as unit generators and event streams, provides valuable insight into designing
expressive musical systems.

However, its focus on low-level synthesis and real-time interaction contrasts with the goals of this thesis,
which emphasize structured composition and deterministic behavior.
Motivo aims to provide a higher-level representation of musical concepts, enabling clearer expression of
temporal relationships and compositional structure, while still allowing integration with synthesis backends.

% Source: paper/chapters/03-related-work/musical-representation/midi.tex

\subsection{MIDI}\label{detail-subsec:midi}

\textbf{MIDI} (Musical Instrument Digital Interface) is a standardized protocol for transmitting musical control and
timing information between electronic instruments and computers in real time~\cite{Moog_1986}.
Rather than encoding audio signals, MIDI represents music as a stream of discrete digital messages describing
performance actions such as note onset, pitch, and control changes.

%% ---------------------------------------------------------------------------------------------------------------------

\subsubsection{Key Features}

MIDI encodes musical information as structured messages consisting of status and data bytes~\cite{Moog_1986}.
A fundamental example is the \textit{Note-On} message, which specifies both a pitch (key number) and a velocity value
corresponding to how quickly a key is pressed~\cite{Moog_1986}.
These values are discretized into integer ranges (0--127), reflecting MIDI’s 7-bit resolution~\cite{Moog_1986}.

Additional message types include \textit{Note-Off}, \textit{Control Change}, \textit{Program Change}, and
\textit{Pitch Bend}, each representing different aspects of performance~\cite{Moog_1986}.
Messages are transmitted serially at a fixed rate of 31.25 kbit/s, which constrains the amount of data that can be sent
in real time~\cite{Moog_1986}.

Furthermore, MIDI organizes communication into 16 logical channels, allowing multiple instruments to be controlled over
a single connection~\cite{Moog_1986}.

%% ---------------------------------------------------------------------------------------------------------------------

\subsubsection{Limitations}

For the purposes of this thesis, the main limitation of MIDI is not utility, but scope.
Its design is rooted in a keyboard-centric model of musical interaction, where music is represented as discrete note
events supplemented by parameterized controls such as velocity and aftertouch~\cite{Moore_1988}.
This makes it difficult to represent continuous musical gestures without approximating them as streams of discrete
messages~\cite{Moore_1988}.

Moore (1988) argues that MIDI conflates musical structure with performance data and lacks the ability to represent
higher-level abstractions such as phrasing, hierarchy, and harmonic relationships~\cite{Moore_1988}.
As a result, MIDI functions primarily as a performance protocol rather than a compositional
representation~\cite{Moore_1988}.

Additionally, the finite resolution of MIDI data (typically 7-bit, or 14-bit for certain controls) introduces
quantization effects, particularly noticeable in pitch and dynamics~\cite{Moore_1988}.

Timing is another critical limitation: because MIDI messages are transmitted serially, dense musical passages can
introduce latency and jitter~\cite{Moog_1986}.
Even under typical conditions, delays on the order of 10--30 milliseconds may occur due to processing and transmission
constraints~\cite{Moog_1986}.

%% ---------------------------------------------------------------------------------------------------------------------

\subsubsection{Relevance to Motivo}

MIDI highlights the strengths and weaknesses of low-level, event-based representations of music.
While it excels at interoperability and real-time performance control, it lacks the expressive power needed to encode
higher-level musical structures.

This distinction is important for the architecture proposed in this thesis.
MIDI remains a valuable backend target for rendering and interchange, but it is not sufficient as the primary
representation for structured composition.
This motivates a domain-specific language that separates compositional intent from performance encoding,
provides richer abstractions for
musical structure, and can ultimately compile to formats such as MIDI for execution.

% Source: paper/chapters/03-related-work/musical-representation/abc_notation.tex

\subsection{ABC Notation}\label{detail-subsec:abc-notation}

\textbf{ABC notation} is a text-based music notation system designed to be comprehensible by both humans and computers,
allowing musical scores to be represented using plain text characters such as letters, digits, and
punctuation~\cite{Walshaw_2011}.

It was developed as a lightweight alternative to traditional staff notation, enabling easy storage, sharing, and
processing of musical compositions in digital form.
The system encodes music declaratively, combining structured metadata with symbolic note representations.
An ABC file typically consists of a \textit{header}, which defines musical properties such as title, meter,
and key, and a \textit{tune body}, which contains the actual sequence of musical events~\cite{Walshaw_2011}.

%% ---------------------------------------------------------------------------------------------------------------------

\subsubsection{Key Features}

A central component of \textbf{ABC notation} is its use of \textit{information fields}, which provide structured
metadata about a musical piece.
These fields, such as \texttt{T:} for title, \texttt{M:} for meter, \texttt{L:} for unit note length, and \texttt{K:}
for key, allow the representation of contextual musical information alongside the score~\cite{Walshaw_2011}.

Musical notes are represented using letters from \texttt{A} to \texttt{G}, with octave modifiers applied through commas
and apostrophes.
Durations are expressed relative to a base unit length and can be modified using numeric multipliers or fractional
notation~\cite{Walshaw_2011}.

The language supports a variety of musical constructs, including chords, rests, tuplets, slurs, and ties.
Rhythmic variations can be expressed using constructs such as broken rhythm markers (\texttt{>} and \texttt{<}),
enabling compact encoding of common rhythmic patterns~\cite{Walshaw_2011}.
ABC notation also provides mechanisms for structuring compositions through bar lines, repeats, and multiple voices.
The \texttt{V:} field enables multi-voice compositions, while additional directives allow for annotations, lyrics, and
formatting instructions~\cite{Walshaw_2011}.

Furthermore, ABC notation is widely supported by external tools capable of converting textual representations into
typeset sheet music or MIDI playback, making it a versatile format for both documentation and performance.

%% ---------------------------------------------------------------------------------------------------------------------

\subsubsection{Limitations}

For the purposes of this thesis, the main limitation of \textbf{ABC notation} is again one of scope.
Despite its expressiveness as a notation system, \textbf{ABC notation} is not designed as a programming language and
lacks support for higher-level abstractions.
It does not provide constructs for modular composition, reusable patterns, or hierarchical structuring beyond basic
mechanisms such as macros and repeated sections.

Additionally, the language does not explicitly model time as a first-class concept, instead relying on implicit ordering
of notes and durations, which can make it difficult to represent complex temporal relationships.
The system also conflates multiple concerns, including musical representation, metadata, and layout directives, which
may reduce clarity when attempting to use it as a foundation for compositional systems.

Finally, ABC notation does not include capabilities for sound synthesis or real-time interaction, relying instead on
external tools for playback and rendering.

%% ---------------------------------------------------------------------------------------------------------------------

\subsubsection{Relevance to Motivo}

\textbf{ABC notation} demonstrates how musical structures can be mapped into a clear and human-readable textual format,
providing valuable insight into the design of accessible music-oriented languages.
Its use of structured metadata and symbolic representations offers a strong foundation for expressing musical ideas in a
compact and interpretable way.

However, \textbf{ABC notation} serves a different role from Motivo proposed in this thesis.
It is effective as a notation and exchange format, but it is not intended to function as an executable, high-level
composition language.
This contrast helps clarify why Motivo must combine readability with abstraction, modularity, and explicit
temporal semantics.

% Source: paper/chapters/03-related-work/summary.tex

\section{Summary and Research Gaps}\label{detail-sec:summary-and-research-gaps}

The systems discussed in this chapter illustrate a wide spectrum of approaches to music programming and representation,
each addressing different aspects of musical creation.
High-level composition systems such as \textbf{Sonic Pi}, \textbf{Tidal}, and \textbf{HMusic} prioritize musical
abstraction, expressiveness, and interactive composition workflows.
However, they often rely on sequential or pattern-based models that can limit explicit control over complex temporal
relationships.

Systems for real-time audio control and synthesis, including \textbf{ChucK} and \textbf{SuperCollider}, provide precise
control over timing and sound synthesis.
While these systems offer high flexibility, they operate at a relatively low level of abstraction, requiring users to
manage details related to signal processing, concurrency, and timing explicitly.

Formats for musical representation and exchange such as \textbf{MIDI} and \textbf{ABC notation} focus on encoding
musical information for storage, transmission, and playback\@.
While widely adopted, these formats serve different goals from a composition-oriented domain-specific
language\@: \textbf{MIDI} is effective
as an interchange and playback format, and \textbf{ABC notation} is effective as a human-readable notation format.
Neither is intended to serve as a high-level, executable model for structured composition.

Table~\ref{detail-tab:related-work-comparison} summarizes the main systems and languages discussed in this chapter and
highlights the design trade-offs that motivate Motivo\@.

\begin{table}[htbp]
  \centering
  \footnotesize
  \renewcommand{\arraystretch}{1.5}
  \begin{tabular}{p{2.2cm}p{2.5cm}p{4cm}p{5cm}}
    \toprule
    \textbf{System / Language} & \textbf{Class} & \textbf{Temporal / structural model} & \textbf{Main gap
    relative to this thesis} \\
    \midrule
    Sonic Pi & High-level composition & Sequential execution with notes, sleeps, and simple patterns &
    Limited structural abstraction and weak support for complex temporal relationships \\
    Tidal & High-level composition & Time-varying patterns with functional composition over discrete or
    continuous time & Haskell-centric syntax and focus on event streaming rather than direct compilation or
    explicit timeline structure \\
    HMusic & High-level composition & Grid-like patterns and parallel tracks for rhythmic structure & Relies
    on Sonic Pi, remains pattern-centric, and does not model higher-level timelines explicitly \\
    ChucK & Real-time audio control and synthesis & Explicit time advancement and synchronized concurrent
    shreds & Low-level time management makes high-level compositional structure harder to express \\
    Super\-Collider & Real-time audio control and synthesis & Object-oriented and functional model with
    UGens, patterns, and streams & Powerful but synthesis-oriented, with substantial complexity and weaker
    support for structured high-level composition \\
    MIDI & Musical representation & Discrete message stream of note and control events & Strong backend and
    interchange format, but too low-level to serve as the primary model for structured composition \\
    ABC notation & Musical representation & Declarative textual notation with implicit ordering, bars, and
    voices & Clear notation format, but not an executable high-level composition language with explicit
    temporal abstraction \\
    \bottomrule
  \end{tabular}
  \caption{Comparison of the main systems and languages discussed in this
  chapter.}\label{detail-tab:related-work-comparison}
\end{table}

From this analysis, several key gaps relative to the goals of this thesis emerge.
First, many systems lack an explicit and flexible model for representing time, often relying on implicit sequencing or
pattern repetition.
Second, there is a gap between low-level audio programming languages and high-level compositional tools, with few
systems providing both expressive power and structured abstraction.
Third, existing approaches often conflate composition with execution, making it difficult to separate musical intent
from performance details.

These gaps motivate the development of the domain-specific language proposed in this thesis, which aims to
provide a timeline-based model for music composition, support deterministic and structured representations of musical
ideas, and enable modular and extensible compilation to multiple output formats.
